\begin{solution}
Let \(P_{ij}=\delta_{ij}-n_in_j\).  Then
\[
\Lambda_{ij,kl}=P_{i(k}P_{l)j}-\frac12P_{ij}P_{kl},
\qquad
S_{ij}^{\rm TT}=\Lambda_{ij,kl}S_{kl}.
\]
Using \(P^2=P\), \(P_{ij}n^j=0\), and \(P^i{}_i=2\) verifies
\(\Lambda^2=\Lambda\), \(n^iS_{ij}^{\rm TT}=0\), and
\(\delta^{ij}S_{ij}^{\rm TT}=0\).
\end{solution}

\begin{solution}
For orbital frequency \(\Omega\) and distance \(D\),
\[
h_+=\frac{4\mathcal M^{5/3}\Omega^{2/3}}{D}
\frac{1+\cos^2\iota}{2}\cos(2\Omega t),
\quad
h_\times=\frac{4\mathcal M^{5/3}\Omega^{2/3}}{D}
\cos\iota\,\sin(2\Omega t).
\]
Face-on, the amplitudes are equal in quadrature, giving circular
polarization.  Edge-on, \(h_\times=0\), giving linear polarization.
\end{solution}

\begin{solution}
TT projection of an axisymmetric quadrupole gives an amplitude proportional
to
\[
h_+\propto\frac{\ddot Q_{zz}}{D}\sin^2\theta,\qquad h_\times=0.
\]
The flux is proportional to \(\dot h_+^2\), hence to
\(\sin^4\theta\).  It vanishes on the symmetry axis and is maximal in the
equatorial plane.
\end{solution}

\begin{solution}
For total mass \(M\) and reduced mass \(\mu\),
\[
E=-\frac{\mu M}{2r},\qquad
P_{\rm GW}=\frac{32}{5}\frac{\mu^2M^3}{r^5}.
\]
Energy balance yields
\[
\dot r=-\frac{64}{5}\frac{\mu M^2}{r^3},\qquad
\frac{\dot P}{P}=\frac32\frac{\dot r}{r},\qquad
t_c=\frac{5r_0^4}{256\mu M^2}.
\]
\end{solution}

\begin{solution}
With \(f=\Omega/\pi\),
\[
\dot f=\frac{96}{5}\pi^{8/3}\mathcal M^{5/3}f^{11/3},
\qquad
t_c-t=\frac{5}{256}(\pi f)^{-8/3}\mathcal M^{-5/3}.
\]
Integrating \(f/\dot f\) gives
\[
N(f_1,f_2)=\frac{1}{32\pi^{8/3}\mathcal M^{5/3}}
\left(f_1^{-5/3}-f_2^{-5/3}\right).
\]
\end{solution}

\begin{solution}
The leading flux is
\[
\frac{\dd J_i}{\dd t}
=-\frac{2}{5}\epsilon_{ijk}
\left\langle\ddot Q_{j\ell}\dddot Q_{k\ell}\right\rangle.
\]
For a circular binary this gives
\(\dot J=-(32/5)\mu^2M^{5/2}r^{-7/2}\).  Since
\(\Omega=\sqrt{M/r^3}\), multiplication by \(\Omega\) reproduces
\(\dot E=-(32/5)\mu^2M^3r^{-5}\).
\end{solution}

\begin{solution}
With the standard one-sided convention,
\[
h_c^2(f)=2fS_h(f),\qquad
\Omega_{\rm GW}(f)
=\frac{1}{\rho_c}\frac{\dd\rho_{\rm GW}}{\dd\ln f}
=\frac{2\pi^2}{3H_0^2}f^2h_c^2(f)
=\frac{4\pi^2}{3H_0^2}f^3S_h(f).
\]
Different spectral conventions move factors of two but not the observable
relations.
\end{solution}

\begin{solution}
The zero-frequency part of the wave equation gives schematically
\[
\Delta h_{ij}^{\rm TT}(\hat{\bm N})
=\frac{4}{D}\left[
\int\frac{n_in_j}{1-\hat{\bm N}\cdot\bm n}
\frac{\dd E}{\dd\Omega_{\bm n}}\,\dd\Omega_{\bm n}
\right]^{\rm TT}.
\]
For initially comoving detector masses separated by \(\xi^j\), integrating
geodesic deviation through the burst yields
\[
\Delta\xi^i=\frac12\Delta h^i{}_j{}^{\rm TT}\xi^j,
\]
a permanent displacement.
\end{solution}
